\section{统计建模：把 benchmark score 当作随机变量}

\subsection{最小数学框架}
令某任务上的观测分数为 $S$，理想能力为 $\mu$，噪声为 $\epsilon$，则
$$
S = \mu + \epsilon,\quad \mathbb{E}[\epsilon]=0,\quad \mathrm{Var}(\epsilon)=\sigma^2.
$$
如果只报告 $S$ 而不报告 $\sigma$，我们无法判断 ``改进'' 是否真实。进一步对多次运行 $S_1,\dots,S_n$ 取均值 $\bar{S}$，其方差缩放为 $\sigma^2/n$，这解释了为什么多次重复评测是必要步骤。

\begin{knowledgebox}{为什么单次跑分不够}
单次跑分只给你一个 realization，无法告诉你 score distribution。没有分布信息，就无法做显著性判断。
\end{knowledgebox}

\subsection{bootstrap 与置信区间}
课程里反复提到 bootstrap 的价值：在样本规模有限、分布未知时，bootstrap 提供了实际可操作的 uncertainty estimate。对于 benchmark reporting，可以使用 percentile CI（如 95\% CI）作为标准报告项。

\begin{lstlisting}[caption={Benchmark 置信区间计算的最小伪代码}]
scores = run_eval_k_times(model, benchmark, k=20)
mean_score = np.mean(scores)
ci_low, ci_high = bootstrap_ci(scores, confidence=0.95)
report(mean_score, ci_low, ci_high)
\end{lstlisting}

\begin{importantbox}{报告规范建议}
每个主结果至少包含：mean、std、95\% CI、重复次数 $k$、协议版本号。
\end{importantbox}

\subsection{噪声可预测性的工程收益}
一旦噪声被量化，决策就会变得更理性。比如模型改进 pipeline 中，可以把 ``是否上线'' 从单分阈值改成 ``显著改进阈值''：只有当新模型相对旧模型在关键子集上显著优于阈值，才进入下一阶段。这会显著减少 ``看起来进步，线上退化'' 的事故。

\begin{warningbox}{未校准噪声会误导训练方向}
如果评测噪声被误当作能力增益，训练团队会被带向错误目标，浪费计算预算和研究周期。
\end{warningbox}

\subsection{本章小结}
这讲最重要的统计思想很朴素：分数必须和不确定性一起报告。bootstrap、重复评测和显著性检验不是锦上添花，而是 benchmark 能被信任的最低门槛。

